\subsubsection{Compito C6: Commentare un file della community}

Il compito C6 richiedeva all’utente di aprire un documento precedentemente individuato all’interno della sezione ``Community'' e inserire un commento testuale specifico (\textit{``Documento molto utile, grazie!''}). L’obiettivo era valutare l’usabilità delle funzionalità di interazione della piattaforma, verificando la facilità di individuazione dell’area commenti e la chiarezza del processo di pubblicazione.

\paragraph{Analisi delle prestazioni generali}

Come si può osservare dai dati aggregati per il compito C6 (Tabella~\ref{tab:6-c6-1}), il tasso di completamento (Success Rate) si è mantenuto al 100,00\% in entrambe le versioni del sistema. La versione post-riprogettazione (B) evidenzia una riduzione del 33,33\% nella media degli errori (da 0,40 a 0,27) e una diminuzione del 9,09\% della difficoltà percepita (da 1,47 a 1,33).

Parallelamente, si registra un lieve incremento della mediana del tempo di esecuzione (+14,29\%, da 0.35.00 a 0.40.00) e della mediana dei click (+16,67\%, da 6 a 7).

\begin{table}[H]
	\centering
	\small
	\renewcommand{\arraystretch}{1.2}
	\caption{Metriche prestazionali globali aggregate per il compito C6.}
	\label{tab:6-c6-1}
	\begin{tabular}{|l|c|c|c|}
		\hline
		\rowcolor{gray!30}
		\textbf{Metrica} & \textbf{Versione originale (A)} & \textbf{Versione post-riprogettazione (B)} & \textbf{Variazione} \\ \hline
		Success Rate     & 100,00\%                        & 100,00\%                                   & 0,00\%              \\ \hline
		Media Errori     & 0,40                            & 0,27                                       & -33,33\%            \\ \hline
		Mediana Tempo    & 0:35                            & 0:40                                       & +14,29\%            \\ \hline
		Mediana Click    & 6                               & 7                                          & +16,67\%            \\ \hline
		Media Difficoltà & 1,47                            & 1,33                                       & -9,09\%             \\ \hline
	\end{tabular}
\end{table}

L'aumento contenuto di tempi e click, accompagnato da una riduzione degli errori e della difficoltà percepita, suggerisce che la versione B introduca un'interazione maggiormente guidata e controllata. La presenza di feedback espliciti e di meccanismi di validazione sembra aver incoraggiato gli utenti a prestare maggiore attenzione durante l'inserimento del commento, riducendo le possibilità di errore a fronte di un lieve incremento delle azioni richieste.

\paragraph{Analisi qualitativa: commenti e osservazioni degli utenti}

Le osservazioni raccolte tramite protocollo \textit{thinking aloud} forniscono una chiave interpretativa utile per comprendere le differenze emerse tra le due versioni.

Per la versione post-riprogettazione (B), i feedback sono stati prevalentemente positivi:

\begin{itemize}
	\item \textbf{Feedback e validazione apprezzati}: diversi utenti hanno notato e apprezzato la presenza di controlli e messaggi di validazione durante l'inserimento del commento. In particolare, è stato esplicitamente menzionato il filtro per contenuti inappropriati (``Mi piace che ci sia un filtro sui contenuti inappropriati''), percepito come un elemento di affidabilità e cura della piattaforma.

	\item \textbf{Maggiore consapevolezza dell'azione}: la presenza di feedback immediati ha contribuito a rendere più chiaro l'esito dell'interazione, trasmettendo agli utenti una maggiore sicurezza durante la pubblicazione del commento.
\end{itemize}

Per la versione originale (A), sono emerse invece alcune criticità riconducibili principalmente alla minore presenza di feedback e alla ridotta visibilità degli elementi interattivi. In diversi casi gli utenti hanno mostrato esitazioni nell'individuazione dell'area dedicata ai commenti e hanno manifestato minore certezza sull'effettiva esecuzione dell'azione richiesta.

\paragraph{Valutazione dell’effetto di apprendimento}

Analizzando i dati in base all'ordine di esecuzione (Tabella~\ref{tab:6-c6-2}), emergono dinamiche di apprendimento tipiche di un task fortemente basato sulla localizzazione di specifici elementi dell'interfaccia.

Nel gruppo A--B, il tempo mediano passa da 0.40 nella versione A a 0.30 nella versione B, con una riduzione del 25,00\%. Contestualmente si osserva una diminuzione degli errori (da 0,67 a 0,22), suggerendo che gli utenti abbiano beneficiato sia dell'esperienza maturata durante la prima esecuzione sia della maggiore chiarezza della versione riprogettata.

Nel gruppo B--A, l'esperienza acquisita inizialmente sulla versione B sembra aver favorito una più rapida individuazione delle funzionalità necessarie anche nella versione originale. La familiarità con la posizione dell'area commenti e con il flusso operativo ha probabilmente attenuato parte delle difficoltà presenti nell'interfaccia A, contribuendo a contenere l'impatto delle sue limitazioni progettuali.

\begin{table}[H]
	\centering
	\small
	\renewcommand{\arraystretch}{1.2}
	\caption{Confronto dei tempi mediani in funzione dell'ordine di esecuzione del test.}
	\label{tab:6-c6-2}
	\begin{tabular}{|>{\raggedright\arraybackslash}p{3cm}|>{\centering\arraybackslash}p{3.8cm}|>{\centering\arraybackslash}p{3.8cm}|>{\centering\arraybackslash}p{2.2cm}|}
		\hline
		\rowcolor{gray!30}
		\textbf{Ordine di Esecuzione} & \textbf{Mediana Tempo Versione A (minuti)} & \textbf{Mediana Tempo Versione B (minuti)} & \textbf{Variazione} \\ \hline
		Ordine A--B                   & 0.40                                       & 0.30                                       & -25,00\%            \\ \hline
		Ordine B--A                   & 0.40                                       & 0.55                                       & +37,50\%            \\ \hline
	\end{tabular}
\end{table}
